\documentclass[10pt,a4paper]{article}

\usepackage[a4paper,left=1.45cm,right=1.45cm,top=1.2cm,bottom=1.2cm]{geometry}
\usepackage[utf8]{inputenc}
\usepackage{setspace}
\usepackage[table]{xcolor}
\usepackage{colortbl}
\usepackage{array}
\usepackage{tabularx}
\usepackage{booktabs}
\usepackage{graphicx}
\usepackage{float}
\usepackage{caption}
\usepackage{ragged2e}

\pagestyle{empty}
\setlength{\parindent}{0pt}
\setlength{\parskip}{0pt}
\setlength{\tabcolsep}{4pt}
\renewcommand{\arraystretch}{1.16}
\captionsetup{font=footnotesize}

\definecolor{ReportBlue}{HTML}{0D2C54}
\definecolor{LightBlue}{HTML}{1E5799}
\definecolor{LightGray}{HTML}{F4F6F9}
\definecolor{MidGray}{HTML}{C8CDD6}
\definecolor{GoodGreen}{HTML}{1E7D46}
\definecolor{MetricRed}{HTML}{C0392B}
\definecolor{CaptionGray}{HTML}{666666}
\definecolor{BodyGray}{HTML}{333333}
\definecolor{FooterGray}{HTML}{999999}
\definecolor{KPIBlueBG}{HTML}{EAF0FB}
\definecolor{KPIGreenBG}{HTML}{E8F5EE}

\newcommand{\SectionBar}[1]{%
  \vspace{0.22em}
  \noindent\colorbox{ReportBlue}{%
    \parbox{\dimexpr\linewidth-2\fboxsep\relax}{%
      \vspace{2.3pt}\textcolor{white}{\bfseries\footnotesize #1}\vspace{2.3pt}
    }%
  }%
  \vspace{0.28em}
}

\newcommand{\ReportTitle}{Assignment IV -- Sparse S\&P 500 Portfolio Replication}
\newcommand{\ReportSubtitle}{DA6701 | Benchmark: \texttt{\string^GSPC} daily log returns | Train: Jan 2020--Dec 2024 | Validation: Jan 2025--Jun 2025 | Holdout: Jul 2025--Dec 2025}

\newcommand{\ReadableFiles}{578}
\newcommand{\FinalStocks}{553}
\newcommand{\TotalExcluded}{31}
\newcommand{\BestMethod}{GA}
\newcommand{\BestTE}{4.63\%}
\newcommand{\BestIR}{-0.10}

\newcommand{\TEFigure}{../figures/te_vs_k.png}
\newcommand{\IRFigure}{../figures/information_ratio.png}
\newcommand{\SectorFigure}{../figures/sector_drift_ga.png}

\newcommand{\FooterText}{Tools: Python 3, Pandas, NumPy, scikit-learn, PyTorch, SciPy (SLSQP) | Methods: Lasso, Autoencoder, Genetic Algorithm | Final universe: 553 stocks + benchmark}

\begin{document}
\color{BodyGray}

\begin{center}
  {\fontsize{15}{17}\selectfont\textbf{\textcolor{ReportBlue}{\ReportTitle}}}\par
  \vspace{1pt}
  {\fontsize{8}{10}\selectfont\textcolor{CaptionGray}{\ReportSubtitle}}
\end{center}
\vspace{3pt}
\noindent\textcolor{ReportBlue}{\rule{\linewidth}{2pt}}
\vspace{5pt}

\noindent
\begin{minipage}[t]{0.496\linewidth}
  \arrayrulecolor{LightBlue}
  \setlength{\tabcolsep}{3pt}
  \begin{tabular}{|>{\centering\arraybackslash}p{0.30\linewidth}|>{\centering\arraybackslash}p{0.30\linewidth}|>{\centering\arraybackslash}p{0.30\linewidth}|}
    \hline
    \cellcolor{KPIBlueBG}{\fontsize{13}{15}\selectfont\textbf{\textcolor{LightBlue}{\ReadableFiles}}} &
    \cellcolor{KPIBlueBG}{\fontsize{13}{15}\selectfont\textbf{\textcolor{GoodGreen}{\FinalStocks}}} &
    \cellcolor{KPIBlueBG}{\fontsize{13}{15}\selectfont\textbf{\textcolor{MetricRed}{\TotalExcluded}}} \\
    \cellcolor{KPIBlueBG}{\fontsize{7}{9}\selectfont\textcolor{CaptionGray}{Readable CSVs}} &
    \cellcolor{KPIBlueBG}{\fontsize{7}{9}\selectfont\textcolor{CaptionGray}{Final Stocks}} &
    \cellcolor{KPIBlueBG}{\fontsize{7}{9}\selectfont\textcolor{CaptionGray}{Excluded Files / Names}} \\
    \hline
  \end{tabular}
\end{minipage}
\hfill
\begin{minipage}[t]{0.496\linewidth}
  \arrayrulecolor{GoodGreen}
  \setlength{\tabcolsep}{3pt}
  \begin{tabular}{|>{\centering\arraybackslash}p{0.30\linewidth}|>{\centering\arraybackslash}p{0.30\linewidth}|>{\centering\arraybackslash}p{0.30\linewidth}|}
    \hline
    \cellcolor{KPIGreenBG}{\fontsize{13}{15}\selectfont\textbf{\textcolor{LightBlue}{\BestMethod}}} &
    \cellcolor{KPIGreenBG}{\fontsize{13}{15}\selectfont\textbf{\textcolor{MetricRed}{\BestTE}}} &
    \cellcolor{KPIGreenBG}{\fontsize{13}{15}\selectfont\textbf{\textcolor{GoodGreen}{\BestIR}}} \\
    \cellcolor{KPIGreenBG}{\fontsize{7}{9}\selectfont\textcolor{CaptionGray}{Best Method}} &
    \cellcolor{KPIGreenBG}{\fontsize{7}{9}\selectfont\textcolor{CaptionGray}{Holdout TE @ k=50}} &
    \cellcolor{KPIGreenBG}{\fontsize{7}{9}\selectfont\textcolor{CaptionGray}{Holdout IR @ k=50}} \\
    \hline
  \end{tabular}
\end{minipage}

\vspace{5pt}

\SectionBar{1 \;|\; Problem Setup, Data Pipeline, and Model Design}
{\fontsize{7.55}{9.55}\selectfont\justifying
The task is to replicate the S\&P 500 daily return profile with a sparse long-only portfolio under the cardinality constraint $k \leq 50$. After validating the raw OHLCV files, we retained \ReadableFiles{} readable stock histories, excluded 6 issue files (5 empty CSVs and one non-price cache file), removed 24 late entrants / early exits, and dropped one additional ticker (CVG) because it still contained 25 missing training prices. The final model-ready matrix covers \texttt{2020-01-03} to \texttt{2025-12-31} and contains 553 stocks plus the benchmark. Validation and holdout price holes were forward-filled before return construction, while the training block was kept hole-free by design.

\vspace{3pt}
The cardinality-constrained tracking problem is an MIQP and is therefore NP-hard. To keep the comparison fair, all methods use the same two-stage pipeline: \textbf{(i)} method-specific stock selection and \textbf{(ii)} a shared long-only quadratic program that minimizes the variance of active returns subject to $\sum_i w_i = 1$ and $w_i \ge 0$. We report annualized Tracking Error (TE) and Information Ratio (IR), where TE should be low and IR should stay close to zero for a true index tracker.
}

\vspace{4pt}
{\fontsize{7.1}{8.9}\selectfont
\arrayrulecolor{MidGray}
\rowcolors{2}{LightGray}{white}
\begin{tabularx}{\linewidth}{|>{\bfseries\RaggedRight\arraybackslash}p{1.85cm}|>{\RaggedRight\arraybackslash}X|>{\RaggedRight\arraybackslash}p{4.2cm}|}
  \hline
  \rowcolor{ReportBlue}
  \textcolor{white}{\textbf{Method}} & \textcolor{white}{\textbf{Selection Stage}} & \textcolor{white}{\textbf{Implementation Note}} \\
  \hline
  Lasso & Non-negative Lasso regression on training benchmark returns; alpha swept on a log-grid to produce sparse portfolios of different $k$. & Fast, interpretable, and directly supervised by the benchmark, but sensitive to regime shocks and concentration. \\
  Autoencoder & Symmetric autoencoder (\texttt{553 $\rightarrow$ 128 $\rightarrow$ 32 $\rightarrow$ 128 $\rightarrow$ 553}); stocks ranked by communality (reconstruction $R^2$) and top-$k$ retained. & Captures shared latent structure, but unsupervised ranking can over-concentrate in sectors with strong common factor loadings. \\
  Genetic Algorithm & Population search over $k$-stock subsets; fitness equals negative training TE, with crossover, mutation, and shared QP reweighting. & Computationally heavier, but best at directly optimizing sparse subset quality under the replication objective. \\
  \hline
\end{tabularx}
}

\vspace{5pt}

\SectionBar{2 \;|\; k = 50 Summary Results}
{\fontsize{7.25}{8.95}\selectfont
\arrayrulecolor{MidGray}
\rowcolors{2}{LightGray}{white}
\begin{tabular}{|>{\bfseries\centering\arraybackslash}p{2.05cm}|>{\centering\arraybackslash}p{1.75cm}|>{\centering\arraybackslash}p{1.55cm}|>{\centering\arraybackslash}p{1.85cm}|>{\centering\arraybackslash}p{1.55cm}|>{\RaggedRight\arraybackslash}p{7.0cm}|}
  \hline
  \rowcolor{ReportBlue}
  \textcolor{white}{\textbf{Method}} &
  \textcolor{white}{\textbf{TE Val}} &
  \textcolor{white}{\textbf{IR Val}} &
  \textcolor{white}{\textbf{TE Holdout}} &
  \textcolor{white}{\textbf{IR Holdout}} &
  \textcolor{white}{\textbf{Interpretation}} \\
  \hline
  Lasso & 5.86\% & -2.29 & 6.78\% & +1.60 & Moderate tracking error, but a large IR sign flip reveals unstable active tilts under the April 2025 tariff shock and subsequent rebound. \\
  Autoencoder & 8.96\% & +1.19 & 8.37\% & -0.74 & Highest TE at $k=50$; the unsupervised ranking produces persistent sector concentration, so the portfolio behaves less like a neutral tracker. \\
  GA & \textbf{4.68\%} & -1.44 & \textbf{4.63\%} & \textbf{-0.10} & Best overall. It delivers the lowest out-of-sample TE on both splits and the holdout IR closest to the ideal target of zero. \\
  \hline
\end{tabular}
}

\vspace{5pt}

\SectionBar{3 \;|\; Required Deliverables: Sparsity vs Tracking Error and Information Ratio}
\noindent
\begin{minipage}[t]{0.59\linewidth}
  \centering
  \includegraphics[width=\linewidth,height=5.9cm,keepaspectratio]{\TEFigure}
\end{minipage}
\hfill
\begin{minipage}[t]{0.39\linewidth}
  \centering
  \includegraphics[width=\linewidth,height=5.9cm,keepaspectratio]{\IRFigure}
\end{minipage}

\vspace{2pt}
\noindent
\begin{minipage}[t]{0.59\linewidth}
  \centering{\fontsize{6.9}{8.4}\selectfont\textcolor{CaptionGray}{Out-of-sample TE falls steeply as $k$ rises from 10 to roughly 50, then flattens. GA dominates the sparse end of the trade-off curve.}}
\end{minipage}
\hfill
\begin{minipage}[t]{0.39\linewidth}
  \centering{\fontsize{6.9}{8.4}\selectfont\textcolor{CaptionGray}{IR quantifies active drift. For a tracker, values near 0 are preferred; GA is closest on holdout.}}
\end{minipage}

\vspace{4pt}
{\fontsize{7.45}{9.3}\selectfont\justifying
The deliverable curve shows the expected diminishing-return pattern: all methods improve rapidly from $k=10$ to around $k=50$, after which TE reductions become incremental. Lasso is especially brittle at small $k$ (TE above 25\% at $k \approx 10$), while GA already remains below 8\% at $k=10$, indicating better subset quality under severe sparsity. The IR plot adds the second required deliverable by highlighting that low TE alone is not sufficient: Lasso and Autoencoder both exhibit unstable directional bets across validation and holdout, whereas GA's holdout IR of -0.10 is the clearest evidence of genuine replication rather than accidental active management.
}

\newpage

\SectionBar{4 \;|\; Required Deliverable: Sector Drift at k = 50}
\noindent
\begin{minipage}[t]{0.62\linewidth}
  \centering
  \includegraphics[width=\linewidth,height=6.9cm,keepaspectratio]{\SectorFigure}
\end{minipage}
\hfill
\begin{minipage}[t]{0.35\linewidth}
  {\fontsize{7.35}{9.1}\selectfont\justifying
  The required sector-drift chart is shown for the GA portfolio at $k=50$, because this is the best overall tracker. Its sector weights remain close to the market-cap weighted S\&P 500 baseline, with only a mild Technology overweight. This makes the portfolio much more defensible as a replication strategy than the alternatives.

  \vspace{4pt}
  Lasso partially satisfies the drift criterion but still overweights Healthcare and underweights Communication Services. Autoencoder fails the sector-neutrality test: it concentrates heavily in Financial Services and Utilities because communality ranking favors stocks explained by the same latent factors.
  }
\end{minipage}

\vspace{5pt}

\SectionBar{5 \;|\; Conclusion}
{\fontsize{7.55}{9.55}\selectfont\justifying
Across all required deliverables, the \textbf{Genetic Algorithm} is the strongest replication method. It achieves the lowest validation and holdout TE at the assignment target of $k=50$, keeps holdout IR closest to zero, and exhibits the cleanest sector alignment with the benchmark. Lasso remains attractive for speed and interpretability but is vulnerable to regime-specific tilts, while the Autoencoder's unsupervised selection objective produces excessive sector concentration. The main empirical takeaway is that, for sparse index replication, directly optimizing subset quality with a search-based method and then reweighting through a common long-only QP yields the most stable out-of-sample tracker.
}

\vspace{7pt}
\noindent\textcolor{MidGray}{\rule{\linewidth}{0.8pt}}
\vspace{2pt}
\begin{center}
  {\fontsize{6.5}{8}\selectfont\textcolor{FooterGray}{\FooterText}}
\end{center}

\end{document}
